\clearpage\begingroup\SourceLeaf{400}\ArithmeticHead{69}{R}{ELEMENTARE UND ANALYTISCHE GEOMETRIE.}
\renewcommand{\thefootnote}{\arabic{footnote})}\setcounter{footnote}{0}
\begin{center}
\textls[100]{Abschnitt IV.}\par\smallskip
\textbf{Elementare und analytische Geometrie.}\par\medskip
\textbf{21.}\par\smallskip
\textls[100]{Allgemeines.}
\end{center}
\noindent\hspace*{\parindent}Im ersten Abschnitt (Nr. 10) ist über verschiedene Untersuchungen von
\textsc{Gauss} berichtet worden, die entweder unmittelbar zur elementaren Geometrie
gehören oder doch eng damit zusammenhängen, bei denen aber das Axiomatische
überwiegt. Auf andere Untersuchungen wurde im dritten Abschnitt
hingewiesen, weil bei ihnen die Anwendung komplexer Grössen mitspielt.
Für deren Gebrauch hatte \textsc{Gauss} eine gewisse Vorliebe, und seine Ausdehnung
erstreckt sich weiter, als man zunächst glauben möchte; \textsc{Gauss} hat sich nämlich
lange Zeit gescheut, mit seiner geometrischen Versinnlichung des Imaginären
öffentlich hervorzutreten, und hat deshalb seine Lösungen in einer davon
befreiten Form dargestellt. Wie gern er mit dem »\(i\)« arbeitet, zeigt übrigens
auch sein Ansatz für das Problem der acht Königinnen, bei dem die Felder
des Schachbrettes mit den Zahlen \(a+ib\) \((a,b=1,2,\ldotp\,\ldotp\,\ldotp,8)\) bezeichnet werden
(Brief an \textsc{Schumacher} vom 27. September 1850, Br. G.-Sch. VI, S. 120).
\par
\textsc{Gauss} hat es erlebt, dass den ursprünglichen, rein geometrischen Überlegungen
und dem später hinzugekommenen Rechnen mit Koordinaten andere
Verfahren zur Lösung geometrischer Aufgaben an die Seite traten, wie der
\textit{Barycentrische Calcul} von \textsc{Möbius} und \textsc{Grassmanns} \textit{Ausdehnungslehre}. Über
den Wert und die Wirksamkeit solcher Methoden hat er sich in dem Brief
an \textsc{Schumacher} vom 15. Mai 1843 mit grosser Klarheit ausgesprochen. Ȇberhaupt
verhält es sich mit allen solchen Kalküls so, dass man durch sie nichts
leisten kann, was nicht auch ohne sie zu leisten wäre; der Vorteil ist aber
der, dass, wenn ein solcher Kalkül dem innersten Wesen vielfach vorkommender
Bedürfnisse korrespondiert, jeder, der sich ihn ganz angeeignet hat, auch ohne
die gleichsam unbewussten Inspirationen des Genies, die niemand \textls[100]{erzwingen}
kann, die dahin gehörigen Aufgaben lösen, ja selbst in so verwickelten Fällen
gleichsam mechanisch lösen kann, wo ohne eine solche Hilfe auch das Genie
ohnmächtig wird. So ist es mit der Erfindung der Buchstabenrechnung
\clearpage\SourceLeaf{401a}\ArithmeticHead{70}{L}{STÄCKEL, GAUSS ALS GEOMETER.}
\noindent überhaupt; so mit der Differentialrechnung gewesen: so ist es auch (wenn auch
in partielleren Sphären) mit \textsc{Lagranges} Variationsrechnung, mit meiner Kongruenzrechnung
und mit \textsc{Möbius}’ Kalkül. Es werden durch solche Konzeptionen
unzählige Aufgaben, die sonst vereinzelt stehen und jedesmal neue
Efforts (kleinere oder grössere) des Erfindungsgeistes erfordern, gleichsam zu
einem organischen Reiche« (W. VIII, S. 298).
\par
Die Arbeiten von \textsc{Gauss}, über die hier berichtet werden soll, betreffen
fast den ganzen Umkreis der elementaren Geometrie, die Anfänge der analytischen
Geometrie eingeschlossen. Eine erste Reihe bezieht sich auf die Eigenschaften
des Dreiecks, des Vierecks und der Vielecke, eine zweite auf den
Kreis und die Kugel, die Kegelschnitte und die Flächen zweiter Ordnung.
Dazu kommen endlich die Beiträge zur sphärischen Trigonometrie, die in
einer Schlussnummer zusammengefasst sind.
\par
Man könnte diesen Teil des Werkes von \textsc{Gauss} übergehen, ohne dass
sein Ruhm geschmälert würde. Allein es gilt dafür das Wort seines Schülers
und Freundes \textsc{Schumacher}: »Deutlich genug ist des Meisters Stempel auch
seinen Erholungen aufgedrückt«\footnote{L. \textsc{Carnot}, \textit{Geometrie der Stellung}, übersetzt von H. C. \textsc{Schumacher}, 2. Teil, Altona 1810, Vorrede, S. II.}.
\clearpage\endgroup
